\section{The domestic clause}
\label{sec:6.5}
The term has a second clause, drawn not from the Charter but from the belligerent's own constitution. A machine takes no part in a use of force by a state whose constitution requires legislative authorization for it, unless that authorization exists. For the United States that means a declaration of war or an authorization for the use of military force that covers the conflict. The War Powers Resolution's sixty days \autocite{wpr1973}, run without one, is evidence the authorization is absent; a concurrent resolution directing withdrawal, adopted by both chambers, is evidence of the same kind, though since Chadha it binds nothing, and a joint resolution passed and vetoed is stronger, because the President had to act to defeat it. The clause has the Charter clause's shape. The term's default is refusal, and what lifts it is an act by a body other than the one that wants the war: a disinterested finding under the Charter, Congress under the Constitution, which gives the war power to the legislature and not to the branch that would fight. The executive's word that it has authorization, or doesn't need it, is its own case. A vendor can hold the clause as it holds the Charter clause, as a term in its license; it becomes standard only by statute, the one route that binds every vendor at once, adopted in peacetime by a legislature with its own reason to give its authorization teeth.\note{
Both clauses read for an act by a body other than the belligerent, a disinterested finding or a legislature's vote, and neither is satisfied by the belligerent describing its own case; what the emergency license takes on the defender's account it takes for the published interval only. A vendor that reads for the act has judged nothing. It has declined to take the interested party's word.
}
Until the statute exists, a vendor holds it as a term, a refusal and a record. Whether a declaration or an authorization exists is a fact about the statute book. Whether an existing authorization covers this conflict, whether the sixty days have run, and whether a ceasefire stops the clock are the executive's claims about its own case, and an old authorization stretched over a new war lifts nothing until Congress says so. It doesn't ask the machine to decide what the Constitution permits, which courts decline to decide. It asks whoever writes the deployment's floor to publish the reading, with the markers it rests on, beside the deployment's translation, where it can be contested, like any other term, through the floor court, the court whose jurisdiction is fixed before the deployment it hears and which hears every claim under the floor (\S\ref{sec:14.3}). The executive's word that its own action is authorized settles nothing. The clause rests on the foundation the rest of the term rests on. A war the legislature never authorized subjects a population abroad to a decision that the state's own people, through their representatives, had no standing to contest either.

\begin{esbox}
\textbf{The record, as of 24 September 2026.} The 2026 campaign supplies every marker the clause reads, and none was met. No declaration was made and no authorization passed; one was introduced in the House on 7 May and never left committee \autocite{hjres176}. The administration reported to Congress under the War Powers Resolution on 2 March, which started the sixty days. On 1 May, the sixtieth day, the President wrote to Congress that the hostilities begun on 28 February had terminated, resting on a ceasefire declared on 7 April, while a naval blockade continued \autocite{cbs2026terminated,lawfare2026sixtydays}. Hostilities resumed in July under no new authorization. In June both chambers adopted a concurrent resolution directing the President to remove forces from hostilities against Iran, and nothing changed \autocite{hconres86,asil2026iranwpr}. The clause doesn't need the votes. What it reads is the absence of an authorization, and by 1 May that absence was on the record in the President's own letter, which conceded that hostilities had begun and claimed only that they had ended. Appendix~\ref{sec:B.6} gives the record vote by vote.
\end{esbox}
